% benefit_society/drafts/zh/manuscript.tex
% P4.2 中文 manuscript 骨架 (plan-r3 衍生论文扩展, 2026-05-15)
% 编译: xelatex + biber + xelatex × 2

% papers_assets/_shared/preamble_zh.tex
% Shared Chinese LaTeX preamble (XeLaTeX + CJKutf8).
% P4.1 (plan-r3 derivative-paper extension, 2026-05-15)
% Compile: xelatex + biber + xelatex × 2

\documentclass[12pt,a4paper]{article}

% --- CJK Fonts (XeLaTeX, cross-platform Q1.1b fallback) ---
\usepackage{fontspec}
\usepackage{xeCJK}
% Windows (local dev): SimSun/SimHei/FangSong
% Linux/macOS (CI): Noto Serif CJK SC / Noto Sans CJK SC (fonts-noto-cjk apt package)
% Last resort: AR PL UMing CN (fonts-arphic-uming) which is preinstalled in many texlive distros
\IfFontExistsTF{SimSun}{%
  \setCJKmainfont{SimSun}[BoldFont=SimHei]%
}{%
  \IfFontExistsTF{Noto Serif CJK SC}{%
    \setCJKmainfont{Noto Serif CJK SC}[BoldFont=Noto Serif CJK SC]%
  }{%
    \setCJKmainfont{AR PL UMing CN}[AutoFakeBold=2]%
  }%
}
\IfFontExistsTF{SimHei}{%
  \setCJKsansfont{SimHei}%
}{%
  \IfFontExistsTF{Noto Sans CJK SC}{%
    \setCJKsansfont{Noto Sans CJK SC}%
  }{%
    \setCJKsansfont{AR PL UKai CN}%
  }%
}
\IfFontExistsTF{FangSong}{%
  \setCJKmonofont{FangSong}%
}{%
  \IfFontExistsTF{Noto Sans Mono CJK SC}{%
    \setCJKmonofont{Noto Sans Mono CJK SC}%
  }{%
    \setCJKmonofont{AR PL UMing CN}%
  }%
}
\setmainfont{Times New Roman}[Ligatures=TeX]
\xeCJKsetup{CJKecglue={\hskip 0.15em plus 0.05em minus 0.05em}}

% --- Math ---
\usepackage{amsmath, amssymb, amsthm}

% --- Tables & Figures ---
\usepackage{booktabs}
\usepackage{graphicx}
\usepackage{longtable}
\usepackage{rotating}
\usepackage{caption}
\captionsetup{font=small,labelfont=bf,skip=4pt}

% --- Layout ---
\usepackage[margin=1in]{geometry}
\usepackage{setspace}
\onehalfspacing
\usepackage{hyperref}
\hypersetup{colorlinks=true,linkcolor=blue,citecolor=blue,urlcolor=blue}

% --- Bibliography ---
\usepackage[backend=biber,style=apa,natbib=true]{biblatex}

% --- Macros ---
\newcommand{\code}[1]{\texttt{#1}}
\usepackage[dvipsnames]{xcolor}
\newcommand{\NoteThis}[1]{\textbf{\textcolor{red}{[TODO: #1]}}}

% --- Theorem environments ---
\newtheorem{prop}{命题}
\newtheorem{hyp}{假设}


\addbibresource{refs.bib}

\title{银行媒体报道中的经营效益态度：\\双极锚点投影法}
\author{\NoteThis{作者姓名}}
\date{\today}

\begin{document}
\maketitle

\begin{abstract}
\NoteThis{最终摘要待定。}
本文基于银行—年度金融科技媒体报道, 采用双极锚点投影 (Jha--Liu--Manela 2026 RFS),
度量中国上市银行的经营效益态度 (\code{benefit\_mean\_full})。我们发现 \NoteThis{核心结论}\@。
\end{abstract}

\section{引言}\label{sec:intro}
\NoteThis{阐释研究问题与贡献。}

\section{文献综述}\label{sec:lit}
三条主线: (1)~双极锚点投影 (Jha--Liu--Manela 2026 \emph{RFS}); (2)~象征--实质采用 (Westphal--Zajac 2001 \emph{ASQ}); (3)~公司金融文本测度 (Hoberg--Phillips 2016 \emph{JPE})。

\section{方法}\label{sec:method}
\subsection{样本与数据}
银行—年度面板覆盖约 164 个银行—年观测 (2017--2018)。
所有数据由 \code{news\_process} 管线 (Phase 9b 基线) 处理,
详见 \code{papers\_assets/V6\_METHODOLOGY\_CHANGES.md}。

\subsection{语调标注 (F1: 三家族 LLM 真三角验证)}
采用三个独立 LLM 家族: \code{gpt-5.4-nano} (OpenAI)、
\code{gemini-3.1-flash-lite} (Google)、\code{deepseek-v4-flash}
(DeepSeek)。OpenAI 与 Google 两家通过 yunwu.ai
(\code{/v1/chat/completions}, OpenAI 兼容) 调用, DeepSeek 直连
\code{api.deepseek.com}; 所有调用 \code{temperature=0}、\code{seed=42},
并使用字节级冻结 prompt (版本 v8)。
家族间一致性以 Cohen's $\kappa$ 与 Krippendorff's $\alpha$ 报告;
多数投票得最终标签, 平局标 \code{undefined} (非 \code{neutral})。

\subsection{篇章级编码 (F2: 去污染)}
三家族全部失败 / 三方分裂者编为\emph{缺失} (非 \code{neutral});
银行--金融科技相关性未通过多数验证者亦视为缺失。
此举消除了 v6 之前 81.7\,\% ``中性'' 标签的污染。

留篇门控仅取全文轨 (\code{llm\_verified\_full}); 实体上下文轨
(\code{ctx}) 与银行—金科同句共现轨 (\code{ss}) 为以全文轨为条件
(conditional-on-full) 的稳健性轨, 其行宇宙是全文轨留篇集的子集。
相较早期"任一轨验证通过即保留"的版本, 更严格的门控会机械性压缩
\code{n\_articles} 及派生份额 (\code{fintech\_share}、
\code{attention\_pct}) --- 这是定义性、方向偏保守的水平位移,
并非数据错误。

\subsection{零观测语义 (F3: NaN 主路径)}
对 \code{n\_articles=0} 的银行—年, 相似度类特征
(\code{auth\_sim\_*}、\code{benefit\_*}、\code{article\_sim\_*})、
比例类 (\code{token\_adopt\_*}、\code{hidden\_adopt\_*}、
\code{subst\_adopt\_*})、\code{fintech\_intensity\_mean} 均编为
\code{NaN}; 计数类保持 0。回归采用列删除 (Allison 2001 / Heckman 1979)。

\subsection{Bartik IV 修剪 (F4)}
约 30 个候选 Bartik 工具变量中, 28 个在 v6 中因违反
Goldsmith-Pinkham \emph{et~al.}\ (2020, \emph{AER}) 外生性被删除。
保留: \code{bartik\_type\_iv\_full} 与 \code{prov\_leave1out\_tone\_full}。

\subsection{嵌入底座 (双主轨、互为稳健性; F5)}
支撑 \code{benefit\_mean\_full} 的双极锚点投影 --- 以及全部嵌入
派生特征 --- 均在来自独立厂商的两个同为主测度的嵌入模型下并行计算、
并排报告、互为稳健性: 阿里 \code{Qwen3-Embedding-8B} (4{,}096 维,
经硅基流动 API; canonical 无后缀列) 与 OpenAI
\code{text-embedding-3-large} (3{,}072 维; \code{\_lg} 后缀列),
遵循严格列对偶 (任一特征两轨齐备或皆无)。各轨进入各自恰好识别的
独立设定, 绝不合并进同一过度识别系统。早期管线版本 canonical 槽为
Google Gemini 嵌入; 自管线 v3.15.0 (2026 年 6 月) 起换为阿里
\code{Qwen3-Embedding-8B}, 全部嵌入派生特征经受控全量重嵌再生,
嵌入空间语义效度以 Spearman 秩相关闸把关 ($\rho \geq 0.80$)。
情感标注 LLM 集成 (F1) 不受此变更影响。

双轨设计并非装饰: 在 Qwen 嵌入空间中, 十对冻结锚点之一 (第 9 对)
的负极投影相对 \code{\_lg} 空间出现翻号 (语义轴锐利度诊断仍通过)。
我们不对冻结锚点做事后调参, 而是保持两条 co-primary 轨按既定口径
并行、对受影响测度做跨厂商双轨并读 --- 这正是跨厂商双轨架构设计
所要吸收的失效模式, 在此如实披露。

\section{数据描述}\label{sec:data}
\NoteThis{Table 1 描述统计参见 \code{papers/benefit_society/tables.py:build\_table1}。}

\section{基准回归}\label{sec:main}
\NoteThis{Table 2 源自 \code{output/tables/table2\_*.csv}; Figure 2 系数图。}

\section{机制检验}\label{sec:mech}
\NoteThis{Table 3 源自 \code{output/tables/table3\_*.csv}; Figure 3。}

\section{异质性}\label{sec:het}
\NoteThis{Table 4 (SOE × 论文专属交互) 源自
\code{output/tables/table4\_heterogeneity.csv}。}

\section{稳健性}\label{sec:robust}
\NoteThis{Table 5 (论文专属稳健性矩阵) 源自
\code{output/tables/table5\_*.csv}。}

% ====================================================================
% V2.0 J-L-M 2026 RFS 稳健性对齐 (Round 9 实施)
% ====================================================================

\subsection{第三方外部指标的外部验证}\label{sec:robust:g4}
\NoteThis{依照 Jha-Liu-Manela (2026, \emph{RFS}) 第3节方法学, 我们将
\texttt{benefit\_mean\_full} 与银行级第三方外部指标进行交叉验证。J-L-M 使用
国家级 credit/GDP/Gini 指标; 我们将\emph{相同}方法学转用到银行级外部 Y:
同业存单 (NCD) 利差、监管处罚频次、SOA 评级。预期方向: NCD 利差 $\downarrow$、
处罚 $\downarrow$、评级 $\uparrow$ (随 benefit-attitude 上升)。Table A9 来自
\code{output/tables/tableA9\_external\_validation.csv}。}

\subsection{安慰剂行业: 银行内部业务线过滤}\label{sec:robust:s1}
\NoteThis{依照 J-L-M (2026) §4.2 placebo industries 方法学, 我们\emph{固定}
主 10 对 fintech 专属锚点方向, 转而变化\emph{文章过滤}范围: 住房贷款、信用卡、
中间业务、网点渠道, 加 1 个中性词基线。如果主信号确为 fintech 领域专属, 投影到
非 fintech 子集上不应预测 ROA。验收阈值: 5 个 placebo 中至少 4 个满足
$|\beta_{\text{placebo}}/\beta_{\text{main}}|<0.3$。Table A10 来自
\code{output/tables/tableA10\_placebo\_business\_lines.csv}。}

\subsection{替代锚点 family: J-L-M canonical 5 对}\label{sec:robust:a2}
\NoteThis{我们使用 J-L-M (2026) Internet Appendix A 中文 5 对 canonical
锚点 (Kozlowski-Taddy-Evans 2019, \emph{ASR} 框架) 重做投影。\emph{这是两个
不同的研究问题, 共享同一方法学}: 我们 10 对锚点针对 fintech 经营效益态度,
J-L-M 5 对针对广义金融情绪。预期 sign-match 一致但 magnitude 会衰减,
因为 J-L-M 广义金融锚点设计上对 fintech 信号不敏感。验收: sign-match 硬要求,
magnitude ratio 仅作 informative 报告。Table A11 来自
\code{output/tables/tableA11\_anchor\_family\_alternative.csv}。}

\section{讨论}\label{sec:discussion}

\subsection{测量与识别的局限性}\label{sec:disc:limit}
\NoteThis{依照 Jha-Liu-Manela (2026) §6.1 议事录传统, 我们公开承认五类局限:
\begin{enumerate}
  \item \textbf{面板时间跨度}: 当前 84 银行、2 年 ($T=2$, 2017--2018) panel
        限制了时间序列检验; fintech 监管事件 event-study 与
        local-projection 脉冲响应推迟到 panel 扩展到 $T\geq10$ (M0.5 里程碑)
        之后。
  \item \textbf{article-level vs sentence-level 投影}: J-L-M 投影 sentence
        embedding (Google Books Ngram corpus 是句子级); 我们的新闻 corpus 是
        文章级。双极投影方法不变, 测量单元有差异。
        \code{core/text/cn\_text\_utils.py} (v3.5.2) 已具备中文 sentence
        splitter 基础设施供后续 sentence-level 稳健性使用。
  \item \textbf{Embedding 模型选择}: 测度在来自独立厂商的两个同为
        主测度的 API 嵌入模型下并行计算 (阿里
        \texttt{Qwen3-Embedding-8B}, 4{,}096 维; OpenAI
        \texttt{text-embedding-3-large}, 3{,}072 维), 并排报告、
        互为稳健性 (早期版本曾使用本地 \texttt{text2vec} 模型与
        Google Gemini 嵌入, 均已退役)。两厂商之外的敏感性已被承认,
        但当前版本未穷尽对比。
  \item \textbf{Cohort 范围}: 分析限于 PAP v1.0-final 84 银行 cohort
        (五大国有 + 股份制 + 城商), 完整 1839 银行 universe 未被分析。
        对非 cohort 机构的推广需要单独的实证检验。
  \item \textbf{跨 corpus 比较}: J-L-M 分析 Google Books Ngram (多语种,
        1870--2009); 我们分析中文金融新闻 (单语种, 2017--2018)。Corpus
        体裁、出版滞后、编辑过滤都存在差异。银行固定效应吸收稳定的
        跨银行测量误差。
\end{enumerate}
}

\subsection{贡献声明}\label{sec:disc:contrib}
\NoteThis{我们相对 Jha-Liu-Manela (2026) 的贡献有三, 这使本文区别于 J-L-M
的"中国版 replication":
\begin{enumerate}
  \item \textbf{观测单元}: 银行年度 firm-level (而非国家年度 macro), 借助
        细粒度固定效应实现 within-bank within-year 识别。
  \item \textbf{outcome 变量}: Token/Hidden 采纳机制
        (Distel-Grimpe-Poetz 2026, Westphal-Zajac 2001 decoupling 框架),
        而非 J-L-M 的宏观 GDP/credit/Gini outcomes。
  \item \textbf{测量设计}: 10 对 fintech-specific 锚点集 + 双轨 full/ctx
        投影 (entity-context masking), 区别于 J-L-M 5 对广义 finance
        canonical 锚点。第\ref{sec:robust:a2}节提供两个锚点 family 的
        交叉验证。
\end{enumerate}
}

\section{结论}\label{sec:concl}
\NoteThis{最终讨论与政策含义。}

\printbibliography

\end{document}
